% !TeX root = forcefields.tex

\documentclass[english,course]{lecture}
\usepackage{amsmath}
\usepackage{hyperref}

% Document metadata
\title{Knowledge Base}
\subtitle{Force Fields: Classical to Machine-Learned}
\shorttitle{Force fields}
\subject{Biophysics}
\author{Alex Payne}
\email{alex.payne@choderalab.org}
\attn{From molecular mechanics to machine learning potentials}
\morelink{http://choderalab.org}
%
% Begin document
\setcounter{tocdepth}{1} % deeper tocs overflow the page and trigger a color-stack error with current TeX Live
\begin{document}

\vfill
\eject

%%%%%%%%%%%%%%%%%%%%%%%%%%%%%%%%%%%%%%%%%%%%%%%%%%%%%%%%%%%%%%%%%%%%%%%%%%%%%%%%%%%%%%%%%%%%%%%%%%%%%%%%%%%%%%%%%%%%%%%%%%%%%%%%%%%%
%% READING
%%%%%%%%%%%%%%%%%%%%%%%%%%%%%%%%%%%%%%%%%%%%%%%%%%%%%%%%%%%%%%%%%%%%%%%%%%%%%%%%%%%%%%%%%%%%%%%%%%%%%%%%%%%%%%%%%%%%%%%%%%%%%%%%%%%%

\section{Recommended reading and supportive texts}

\subsection*{Recommended reading}

\noindent {\bf Understanding Molecular Simulation: From Algorithms to Applications}, Frenkel and Smit \\
Standard reference for how potentials are used in simulation.

\noindent {\bf Geometric Deep Learning: Grids, Groups, Graphs, Geodesics, and Gauges}, Bronstein, Bruna, Cohen, and Veli\v{c}kovi\'{c} \\
Background on symmetry and equivariance in neural networks, which underlies modern ML potentials.
\href{https://arxiv.org/abs/2104.13478}{[arXiv]}

\subsection*{Key papers (to be filled in)}

% TODO: add citations for each family below (Behler--Parrinello, ANI, SchNet, NequIP, MACE, Allegro, foundation models, espaloma).

\vfill
\eject

%%%%%%%%%%%%%%%%%%%%%%%%%%%%%%%%%%%%%%%%%%%%%%%%%%%%%%%%%%%%%%%%%%%%%%%%%%%%%%%%%%%%%%%%%%%%%%%%%%%%%%%%%%%%%%%%%%%%%%%%%%%%%%%%%%%%
%% INTRODUCTION
%%%%%%%%%%%%%%%%%%%%%%%%%%%%%%%%%%%%%%%%%%%%%%%%%%%%%%%%%%%%%%%%%%%%%%%%%%%%%%%%%%%%%%%%%%%%%%%%%%%%%%%%%%%%%%%%%%%%%%%%%%%%%%%%%%%%

\section{What is a force field?}

A force field is a model of the potential energy surface $V(\mathbf{r})$ of a molecular system as a function of nuclear coordinates $\mathbf{r}$. Forces on atoms are the negative gradient,
\[
\mathbf{F}_i = -\nabla_{\mathbf{r}_i} V(\mathbf{r}),
\]
and drive the dynamics described in the molecular dynamics notes. The quality of any simulation---conformational ensembles, binding free energies, kinetics---is limited by the accuracy of $V$, so the choice of force field is one of the most consequential modeling decisions.

The central trade-off is between \emph{accuracy} (how closely $V$ matches quantum mechanics or experiment) and \emph{cost} (energy and force evaluations are performed billions of times in a simulation).

\subsection{How a force field is used}

In an MD simulation (see the molecular dynamics notes), the force field is called at every timestep: given coordinates, return the energy and the force on every atom; the integrator then uses those forces to advance positions and velocities. Because this happens $10^6$--$10^9$ times in a typical simulation, evaluation speed matters as much as accuracy.

% TODO: add a simple worked example (e.g.\ a diatomic or water dimer potential energy curve) and a figure of the potential energy surface.

\subsection{Where force fields sit between QM and MM}

Quantum mechanics (QM) is the most accurate description of the potential energy surface but is expensive; molecular mechanics (MM) force fields are cheap but approximate. Machine-learned force fields are, in effect, a bridge between the two: trained on QM data, evaluated at a cost closer to MM.

\subsection{Goals when improving a force field}

\begin{itemize}
    \item \textbf{Increase speed}: cheaper energy and force evaluations allow longer and larger simulations.
    \item \textbf{Increase accuracy}: closer agreement with QM or experiment (e.g.\ binding free energies).
    \item \textbf{Increase domain of applicability}: cover more chemistry (elements, charge states, protein modifications, drug-like space) and more environments reliably.
\end{itemize}
These goals typically trade off against one another, and much of the research in this area is about moving the Pareto frontier.

\subsection{Properties a good potential should have}

\begin{itemize}
    \item \textbf{Smooth and differentiable}, with forces that are consistent with the energy (conservative forces).
    \item \textbf{Invariant} to translation, rotation, and permutation of identical atoms.
    \item \textbf{Extensive} (or local), so cost scales roughly linearly with system size.
    \item \textbf{Transferable} to molecules and environments not seen during parameterization.
\end{itemize}

\vfill
\eject

%%%%%%%%%%%%%%%%%%%%%%%%%%%%%%%%%%%%%%%%%%%%%%%%%%%%%%%%%%%%%%%%%%%%%%%%%%%%%%%%%%%%%%%%%%%%%%%%%%%%%%%%%%%%%%%%%%%%%%%%%%%%%%%%%%%%
%% CLASSICAL
%%%%%%%%%%%%%%%%%%%%%%%%%%%%%%%%%%%%%%%%%%%%%%%%%%%%%%%%%%%%%%%%%%%%%%%%%%%%%%%%%%%%%%%%%%%%%%%%%%%%%%%%%%%%%%%%%%%%%%%%%%%%%%%%%%%%

\section{Classical (molecular mechanics) force fields}

\subsection{Functional form}

Biomolecular force fields such as AMBER, CHARMM, OPLS, and OpenFF decompose the energy into bonded and nonbonded terms:
\[
V = \sum_{\text{bonds}} \tfrac{1}{2}k_b (b - b_0)^2
  + \sum_{\text{angles}} \tfrac{1}{2}k_\theta (\theta - \theta_0)^2
  + \sum_{\text{torsions}} \sum_n k_n \left[1 + \cos(n\phi - \delta_n)\right]
  + \sum_{i<j} \left[ 4\epsilon_{ij}\left(\frac{\sigma_{ij}^{12}}{r_{ij}^{12}} - \frac{\sigma_{ij}^{6}}{r_{ij}^{6}}\right) + \frac{q_i q_j}{4\pi\varepsilon_0 r_{ij}} \right].
\]
The last term combines Lennard-Jones and Coulomb interactions; the Lennard-Jones potential is introduced in the statistical mechanics notes.

\subsection{Parameterization and atom typing}

Parameters are assigned by matching atoms to \emph{atom types} or chemical environments, then fit to quantum chemical data and experimental observables (densities, enthalpies of vaporization, hydration free energies). Charges are commonly assigned by RESP or AM1-BCC.

\subsection{Limitations}

\begin{itemize}
    \item Fixed functional form: no bond breaking, limited polarization, no charge transfer.
    \item Fixed (non-polarizable) charges in most widely used models.
    \item Parameter coverage for drug-like chemical space is a continual maintenance burden.
\end{itemize}

\vfill
\eject

%%%%%%%%%%%%%%%%%%%%%%%%%%%%%%%%%%%%%%%%%%%%%%%%%%%%%%%%%%%%%%%%%%%%%%%%%%%%%%%%%%%%%%%%%%%%%%%%%%%%%%%%%%%%%%%%%%%%%%%%%%%%%%%%%%%%
%% ML FORCE FIELDS
%%%%%%%%%%%%%%%%%%%%%%%%%%%%%%%%%%%%%%%%%%%%%%%%%%%%%%%%%%%%%%%%%%%%%%%%%%%%%%%%%%%%%%%%%%%%%%%%%%%%%%%%%%%%%%%%%%%%%%%%%%%%%%%%%%%%

\section{Machine learning force fields}

\subsection{Motivation}

Quantum mechanical (QM) methods are accurate but scale steeply with system size. Machine learning (ML) potentials learn a mapping from atomic configuration to energy from QM reference data, aiming for near-QM accuracy at a cost closer to classical force fields.

\subsection{Using ML to improve binding affinity predictions}

A central question for the lab is how machine learning can improve binding free energy predictions. There are three broad routes:

\begin{enumerate}
    \item \textbf{ML--MM force fields}: a model predicts molecular mechanics parameters from the molecular graph (e.g.\ espaloma), keeping the speed and interpretability of the classical functional form.
    \item \textbf{ML potentials}: a neural network replaces the entire energy function (ANI, SchNet, NequIP, MACE, \ldots), either for the whole system or for a region such as the ligand.
    \item \textbf{ML corrections to affinity predictions}: a classical calculation is run and an ML model learns a correction to the resulting free energy or to the energy surface (e.g.\ MM$\to$ML reweighting or $\Delta$-learning).
\end{enumerate}

% TODO: add pros/cons of each route, and example papers.

\subsection{Architectures}

\subsubsection*{Atom-centered descriptors}
Behler--Parrinello networks and ANI-style models write the energy as a sum of atomic contributions, each a neural network function of symmetry-function descriptors of the local environment.

\subsubsection*{Message-passing neural networks}
Atoms are nodes in a graph with edges between neighbors within a cutoff; learned messages are exchanged over several layers (e.g.\ SchNet).

\subsubsection*{Equivariant networks}
Features carry irreducible representations of $SO(3)$ so that vector and tensor quantities rotate correctly with the molecule (NequIP, Allegro, MACE). These typically achieve better data efficiency and accuracy than invariant models.

\subsection{Training}

Training minimizes a loss on energies and, importantly, forces:
\[
\mathcal{L} = \lambda_E \left(E_{\text{pred}} - E_{\text{ref}}\right)^2 + \frac{\lambda_F}{3N}\sum_{i=1}^{N} \left\lVert \mathbf{F}_{i}^{\text{pred}} - \mathbf{F}_{i}^{\text{ref}} \right\rVert^2,
\]
with predicted forces obtained by automatic differentiation of the predicted energy.

\subsection{Datasets and levels of theory}

% TODO: summarize common datasets (ANI-1x, SPICE, QM9, MD17, OMol/OMat-type large-scale sets), the DFT functionals and basis sets used, and the implications for transferability.

\subsection{Pretrained and foundation models}

% TODO: survey of universal/foundation potentials, fine-tuning workflows, and benchmarks.

\subsection{Known challenges}

\begin{itemize}
    \item \textbf{Long-range interactions}: electrostatics and dispersion beyond the cutoff are often missing or handled separately.
    \item \textbf{Charge and spin states}: many models are trained for neutral closed-shell systems only.
    \item \textbf{Extrapolation and robustness}: unphysical behavior in regions of configuration space absent from training data.
    \item \textbf{Cost}: still orders of magnitude slower than classical force fields for large biomolecular systems.
    \item \textbf{Energy conservation and stability} in long MD trajectories.
\end{itemize}

\vfill
\eject

%%%%%%%%%%%%%%%%%%%%%%%%%%%%%%%%%%%%%%%%%%%%%%%%%%%%%%%%%%%%%%%%%%%%%%%%%%%%%%%%%%%%%%%%%%%%%%%%%%%%%%%%%%%%%%%%%%%%%%%%%%%%%%%%%%%%
%% BEST PRACTICES
%%%%%%%%%%%%%%%%%%%%%%%%%%%%%%%%%%%%%%%%%%%%%%%%%%%%%%%%%%%%%%%%%%%%%%%%%%%%%%%%%%%%%%%%%%%%%%%%%%%%%%%%%%%%%%%%%%%%%%%%%%%%%%%%%%%%

\section{Best practices for force field development}

The Open Force Field (OpenFF) community has developed practices for building and validating force fields in a reproducible way.

% TODO: summarize OpenFF practices, e.g. versioned and openly available training and benchmark datasets, automated fitting pipelines (such as the OpenFF BespokeFit/Evaluator/Fitting tooling), held-out benchmark sets, transparent release notes, and reporting of uncertainty. Verify each point against OpenFF documentation before writing.

\vfill
\eject

%%%%%%%%%%%%%%%%%%%%%%%%%%%%%%%%%%%%%%%%%%%%%%%%%%%%%%%%%%%%%%%%%%%%%%%%%%%%%%%%%%%%%%%%%%%%%%%%%%%%%%%%%%%%%%%%%%%%%%%%%%%%%%%%%%%%
%% MD WITH ML POTENTIALS
%%%%%%%%%%%%%%%%%%%%%%%%%%%%%%%%%%%%%%%%%%%%%%%%%%%%%%%%%%%%%%%%%%%%%%%%%%%%%%%%%%%%%%%%%%%%%%%%%%%%%%%%%%%%%%%%%%%%%%%%%%%%%%%%%%%%

\section{Running molecular dynamics with ML potentials}

\subsection{Mixed ML/MM simulations}

Embedding an ML-described region (e.g.\ a ligand) in a classical environment (protein and solvent) balances accuracy and cost. Care is needed for the interface and for the coupling of electrostatics.

\subsection{Free energy calculations}

Alchemical and path-based free energy methods require the potential to be evaluated along a transformation; see the MD and binding-models notes. Using an ML potential for the end states or as a correction to a classical calculation is an active research area.

\subsection{Software tools}

% TODO: list and describe tooling (OpenMM, OpenMM-ML, TorchANI, MACE, ASE, openmmtools, openff toolkit).

\vfill
\eject

%%%%%%%%%%%%%%%%%%%%%%%%%%%%%%%%%%%%%%%%%%%%%%%%%%%%%%%%%%%%%%%%%%%%%%%%%%%%%%%%%%%%%%%%%%%%%%%%%%%%%%%%%%%%%%%%%%%%%%%%%%%%%%%%%%%%
%% CONCLUSION
%%%%%%%%%%%%%%%%%%%%%%%%%%%%%%%%%%%%%%%%%%%%%%%%%%%%%%%%%%%%%%%%%%%%%%%%%%%%%%%%%%%%%%%%%%%%%%%%%%%%%%%%%%%%%%%%%%%%%%%%%%%%%%%%%%%%
\section{Conclusion}

\subsection{Key Takeaways}

\begin{itemize}
    \item A force field defines the potential energy surface; its accuracy bounds the accuracy of any simulation.
    \item Classical force fields are fast and interpretable but limited by their fixed functional form.
    \item ML potentials learn the energy surface from QM data and can approach QM accuracy at far lower cost, subject to training data coverage.
    \item Symmetry-respecting (equivariant) architectures and force-based training are central to modern ML potentials.
    \item Long-range physics, robustness, and cost remain open challenges.
    \item Improvements should be judged against the goals of speed, accuracy, and domain of applicability, using community benchmarks.
\end{itemize}

\end{document}
